\SourceLeaf{328}%
\par\noindent\hspace*{\parindent} Wenn man erwägt, dass in der obigen Entwicklung in Beziehung auf \(Z\) nichts weiter vorausgesetzt ist, als dass es eine von \(u\) abhängige veränderliche Grösse ist, deren Differentialquotient \(\frac{\mathrm{d}Z}{\mathrm{d}u}=z\) durch irgend eine Function von \(X\) und \(Y\) ausgedrückt werde, so kann man die gefundenen Resultate auch unmittelbar auf jede andere in gleichem Falle sich befindende veränderliche Grösse, namentlich auf \(X\) oder \(Y\) selbst, übertragen, wenn man nur anstatt \(L,\mathrm{L},l\), und der verschieden accentuirten \(\mathrm{h}\) beziehungsweise \(T,\mathrm{T},t\) und die verschiedenen \(\mathrm{f}\), oder \(B,\mathrm{B},b\), und die verschiedenen \(\mathrm{g}\) einschiebt. Zunächst gibt uns demnach die Gleichung (4), von welcher hier sonst kein directer Gebrauch gemacht wird, folgende beiden, gleichfalls bis zur vierten Ordnung ausschl. genauen:
\[
\begin{aligned}
T&=\mathrm{T}+\tfrac18(\mathrm{ff}'+\mathrm{f}''\mathrm{g})rr\\
B&=\mathrm{B}+\tfrac18(\mathrm{fg}'+\mathrm{gg}'')rr
\end{aligned}
\]
Man schliesst hieraus zuvörderst, dass \(h\) und \(\mathrm{h}\), als die Werthe von \(z\), je nachdem man \(T\) und \(B\), oder \(\mathrm{T}\) und \(\mathrm{B}\) für \(X\) und \(Y\) substituirt, von einander um eine Grösse zweiter Ordnung verschieden sind, und zwar wird dieser Unterschied, bis auf die vierte Ordnung ausschl. genau, bestimmt durch die Formel
\[
h=\mathrm{h}+\tfrac18(\mathrm{ff}'+\mathrm{f}''\mathrm{g})rr\mathord{.}(\tfrac{\mathrm{d}z}{\mathrm{d}X})+\tfrac18(\mathrm{fg}'+\mathrm{gg}'')rr(\tfrac{\mathrm{d}z}{\mathrm{d}Y})
\]
wo für die partiellen Differentialquotienten \((\frac{\mathrm{d}z}{\mathrm{d}X})\) und \((\frac{\mathrm{d}z}{\mathrm{d}Y})\), oder \(z',z''\) ihre bestimmten Werthe bei \(X=\mathrm{T},\ Y=\mathrm{B}\) anzunehmen sind, nemlich \(\mathrm{h}'\) und \(\mathrm{h}''\). Es ist also, bis auf die vierte Ordnung genau,
\[
\mathrm{h}=h-\tfrac18(\mathrm{ff'h}'+\mathrm{fg'h}''+\mathrm{f''gh}'+\mathrm{gg''h}'')rr
\]
und vermöge der Substitution dieses Werths in der Gleichung (5) wird bis auf die fünfte Ordnung ausschl. genau
\[
\begin{aligned}
l={}&-hr+\tfrac1{24}(2\mathrm{ff'h}'+2\mathrm{fg'h}''+2\mathrm{f''gh}'+2\mathrm{gg''h}''\\
&\qquad-\mathrm{ffh}'''-2\mathrm{fgh}''''-\mathrm{ggh}^{\mathrm{v}})r^3
\end{aligned}
\]

\indent Aus gleichen Gründen wie \(\mathrm{h}\) von \(h\), werden auch \(\mathrm{f},\mathrm{f}',\mathrm{f}''\) u.\,s.\,w. \(\mathrm{g},\mathrm{g}',\mathrm{g}''\) u.\,s.\,w. \(\mathrm{h}',\mathrm{h}''\) u.\,s.\,w. von \(f,f',f''\) u.\,s.\,w. \(g,g',g''\) u.\,s.\,w. \(h',h''\) u.\,s.\,w. beziehungsweise um Grössen zweiter Ordnung verschieden sein, und man kann daher in den eben gegebenen Ausdruck für \(l\) anstatt jener Grössen die letztern ohne Verminderung des Grades der Genauigkeit substituiren. Es ist also gleichfalls bis auf die fünfte Ordnung ausschl. genau
